%!TeX root = ../qualificacao.tex
%!TeX encoding = utf-8
\chapter{Resultados}
\label{ch:results}

Este capítulo reporta a grade completa de 162 treinos. Uma decisão metodológica
precede os números: a razão de retenção $R_c$ não é a métrica principal, porque
nesta grade ela é anticorrelacionada com a acurácia limpa e premia modelos fracos
(Seção~\ref{sec:res:confound}). As conclusões repousam sobre a acurácia absoluta
sob deslocamento, lida contra o acaso e a classe majoritária.

\section{Desempenho em distribuição}
\label{sec:res:clean}

\begin{table}[htb]
  \centering
  \caption{Acurácia de teste limpa (\%), média $\pm$ desvio padrão sobre três
  sementes.}
  \label{tab:res-clean}
  \footnotesize
  \setlength{\tabcolsep}{4pt}
  \begin{tabular}{llccc@{\hspace{12pt}}ccc}
    \toprule
    \multirow{2}{*}{Modelo} & \multirow{2}{*}{Aug.}
      & \multicolumn{3}{c}{CIFAR-10, fração $f$} & \multicolumn{3}{c}{D-Fire, fração $f$} \\
    \cmidrule(lr){3-5}\cmidrule(lr){6-8}
     & & $0{,}33$ & $0{,}66$ & $1{,}00$ & $0{,}33$ & $0{,}66$ & $1{,}00$ \\
    \midrule
    \multirow{3}{*}{E-CapsNet}
      & nenhuma & 63,8$\pm$0,1 & 70,3$\pm$0,4 & 74,2$\pm$0,1 & 87,3$\pm$0,5 & 89,2$\pm$0,2 & 90,2$\pm$0,5 \\
      & padrão  & 78,0$\pm$0,5 & 82,2$\pm$0,1 & 84,8$\pm$0,4 & \textbf{89,0$\pm$0,5} & 92,8$\pm$0,4 & 93,6$\pm$0,4 \\
      & forte   & 76,6$\pm$0,3 & 80,7$\pm$0,1 & 83,0$\pm$0,3 & 86,4$\pm$0,0 & 90,3$\pm$0,4 & 91,9$\pm$0,3 \\
    \addlinespace
    \multirow{3}{*}{ResNet-18}
      & nenhuma & 66,2$\pm$1,9 & 75,3$\pm$0,9 & 78,8$\pm$1,5 & 84,4$\pm$1,7 & 88,7$\pm$0,6 & 91,2$\pm$0,4 \\
      & padrão  & 83,3$\pm$0,5 & 88,6$\pm$0,4 & 90,8$\pm$0,1 & 88,8$\pm$0,9 & 92,7$\pm$0,3 & 94,1$\pm$0,5 \\
      & forte   & \textbf{86,8$\pm$0,6} & \textbf{91,5$\pm$0,1} & \textbf{93,3$\pm$0,1} & \textbf{90,0$\pm$0,7} & \textbf{93,8$\pm$0,2} & \textbf{94,7$\pm$0,2} \\
    \addlinespace
    \multirow{3}{*}{DeiT-Tiny}
      & nenhuma & 52,8$\pm$0,7 & 58,8$\pm$0,5 & 60,9$\pm$2,2 & 82,2$\pm$0,7 & 84,0$\pm$2,5 & 84,6$\pm$2,9 \\
      & padrão  & 59,2$\pm$1,3 & 61,6$\pm$6,6 & 70,2$\pm$3,5 & 76,0$\pm$7,4 & 78,9$\pm$8,8 & 81,0$\pm$5,4 \\
      & forte   & 49,6$\pm$4,0 & 57,8$\pm$8,0 & 68,7$\pm$10,4 & 73,9$\pm$2,0 & 75,1$\pm$5,3 & 77,5$\pm$8,1 \\
    \bottomrule
  \end{tabular}
\end{table}

No CIFAR-10 a ResNet-18 vence nas nove combinações, com margem de $2{,}3$ a
$10{,}8$~pp sobre o Efficient-CapsNet; o melhor resultado capsular ($84{,}8\%$)
fica $8{,}5$~pp abaixo do melhor convolucional ($93{,}3\%$). O DeiT-Tiny,
treinado do zero em imagens pequenas, é o pior dos três. No D-Fire, tarefa
binária e mais fácil, o Efficient-CapsNet fica a $1{,}1$~pp da ResNet-18 na
melhor configuração de cada um e tem a maior média em quatro células com poucos
dados, só uma delas por margem apreciável ($+2{,}9$~pp). A ResNet-18 é a única
arquitetura que melhora ao passar de augmentação padrão para forte
(Tabela~\ref{tab:res-clean}).

\section{Robustez sob transformações não vistas}
\label{sec:res:robust}

\begin{table}[htb]
  \centering
  \caption{Desempenho sob deslocamento em $f = 1{,}00$, sem augmentação em
  teste: acurácia absoluta $A^{c}$, retenção $R_c$ e habilidade retida $S_c$ (\%).}
  \label{tab:res-robust}
  \footnotesize
  \begin{tabular}{llcccccc}
    \toprule
    \multirow{2}{*}{Modelo} & \multirow{2}{*}{Aug.}
      & \multicolumn{3}{c}{Rotação grande} & \multicolumn{3}{c}{Todas as transformações} \\
    \cmidrule(lr){3-5}\cmidrule(lr){6-8}
     & & $A^{c}$ & $R_c$ & $S_c$ & $A^{c}$ & $R_c$ & $S_c$ \\
    \midrule
    \multicolumn{8}{l}{\textit{CIFAR-10}} \\
    \multirow{3}{*}{E-CapsNet}
      & nenhuma & 38,1 & 51,3 & 43,7 & 12,6 & 17,0 & 4,1 \\
      & padrão  & 52,9 & 62,4 & 57,4 & 12,8 & 15,1 & 3,8 \\
      & forte   & 51,4 & 61,9 & 56,6 & 13,8 & 16,7 & 5,2 \\
    \multirow{3}{*}{ResNet-18}
      & nenhuma & 36,4 & 46,2 & 38,4 & 11,6 & 14,7 & 2,3 \\
      & padrão  & 57,8 & 63,6 & 59,1 & 12,3 & 13,6 & 2,9 \\
      & forte   & \textbf{60,0} & 64,4 & 60,1 & 13,4 & 14,4 & 4,1 \\
    \multirow{3}{*}{DeiT-Tiny}
      & nenhuma & 40,0 & 65,8 & 59,0 & 16,4 & 26,9 & 12,6 \\
      & padrão  & 52,8 & 75,1 & 71,0 & 15,9 & 22,6 & 9,8 \\
      & forte   & 52,0 & \textbf{75,7} & \textbf{71,6} & \textbf{16,8} & 24,5 & 11,6 \\
    \midrule
    \multicolumn{8}{l}{\textit{D-Fire}} \\
    \multirow{3}{*}{E-CapsNet}
      & nenhuma & 61,4 & 68,1 & 28,4 & 53,1 & 58,9 & 7,7 \\
      & padrão  & 77,6 & 82,9 & 63,2 & 56,1 & 59,9 & 13,9 \\
      & forte   & 80,0 & 87,0 & 71,5 & 60,5 & 65,8 & 25,0 \\
    \multirow{3}{*}{ResNet-18}
      & nenhuma & 63,4 & 69,5 & 32,4 & 52,5 & 57,5 & 6,0 \\
      & padrão  & 78,7 & 83,7 & 65,2 & 54,9 & 58,4 & 11,2 \\
      & forte   & \textbf{80,9} & 85,4 & 69,1 & \textbf{60,7} & 64,1 & 23,9 \\
    \multirow{3}{*}{DeiT-Tiny}
      & nenhuma & 63,5 & 75,0 & 38,9 & 52,9 & 62,5 & 8,4 \\
      & padrão  & 78,4 & 96,9 & 91,8 & 53,7 & 66,4 & 12,1 \\
      & forte   & 76,0 & \textbf{98,0} & \textbf{94,4} & 59,9 & \textbf{77,2} & \textbf{35,9} \\
    \bottomrule
  \end{tabular}
\end{table}

\paragraph{Rotação.}
A rotação é o teste direto de H1, e o resultado é desfavorável à hipótese. No
CIFAR-10 com augmentação padrão ou forte, a ResNet-18 tem a maior acurácia
absoluta sob rotação ($57{,}8\%$ e $60{,}0\%$ contra $52{,}9\%$ e $51{,}4\%$), e
essa vantagem de até $8{,}6$~pp supera qualquer vantagem capsular da grade. Só
sem augmentação o Efficient-CapsNet fica à frente, por pouco ($38{,}1\%$ contra
$36{,}4\%$). No teste agregado por semente, a diferença de retenção contra a
ResNet-18 é de $+1{,}1$~pp no CIFAR-10 (IC 95\% de $-0{,}3$ a $+2{,}4$;
$p_{\text{Holm}} = 0{,}22$) e de $-1{,}5$~pp no D-Fire ($p_{\text{Holm}} =
0{,}065$). Contra o DeiT-Tiny a retenção capsular é significativamente menor
($-15{,}1$ e $-11{,}7$~pp).

\paragraph{Composição de todas as transformações.}
Sob as seis corrupções combinadas, os modelos perdem quase toda a capacidade de
discriminar: no CIFAR-10 a acurácia fica entre $11{,}1\%$ e $19{,}1\%$ nas 54
células, e no D-Fire 16 das 54 células ficam abaixo da classe majoritária
($53{,}4\%$). Nessa condição, diferenças de $R_c$ comparam padrões de colapso,
não robustez, e por isso não entram nas conclusões.

\subsection{Por que a retenção não é a métrica principal}
\label{sec:res:confound}

O DeiT-Tiny tem a maior retenção em quase toda a Tabela~\ref{tab:res-robust} e
a menor acurácia limpa da Tabela~\ref{tab:res-clean}. Não é coincidência: no
CIFAR-10, a correlação entre acurácia limpa e retenção é de $r = -0{,}44$ sob
rotação e $r = -0{,}91$ sob todas as transformações
(Figura~\ref{fig:res-confound}). O mecanismo é aritmético. A acurácia sob
deslocamento tem piso no acaso; um modelo que cai de $60\%$ para $40\%$ retém
$67\%$, e um que cai de $93\%$ para $60\%$ retém $64\%$, embora esteja $20$~pp
à frente onde importa.

\begin{figure}[htb]
  \centering
  \includegraphics[width=0.75\textwidth]{fig03_retention_confound_cifar_10.png}
  \caption{CIFAR-10: razão de retenção contra acurácia limpa. A relação
  negativa mostra que a retenção, isolada, recompensa modelos fracos.}
  \label{fig:res-confound}
\end{figure}

\section{Eficiência de dados e estabilidade}
\label{sec:res:data}
\label{sec:res:seeds}

\begin{table}[htb]
  \centering
  \caption{Queda de acurácia limpa $\Delta_f$ entre $f = 1{,}00$ e $f = 0{,}33$
  e amplitude média entre sementes (máximo menos mínimo), ambas em pp.}
  \label{tab:res-dataeff}
  \footnotesize
  \begin{tabular}{llcccc@{\hspace{14pt}}cccc}
    \toprule
    \multirow{2}{*}{Conj.} & \multirow{2}{*}{Modelo}
      & \multicolumn{4}{c}{$\Delta_f$} & \multicolumn{4}{c}{Amplitude entre sementes} \\
    \cmidrule(lr){3-6}\cmidrule(lr){7-10}
     & & nenh. & padr. & forte & média & nenh. & padr. & forte & média \\
    \midrule
    \multirow{3}{*}{CIFAR}
      & E-CapsNet & 10,4 & 6,7  & 6,5  & \textbf{7,9} & 0,93 & 0,93 & 0,86 & \textbf{0,91} \\
      & ResNet-18 & 12,6 & 7,5  & 6,5  & 8,9          & 2,01 & 1,17 & 1,06 & 1,42 \\
      & DeiT-Tiny & 8,1  & 11,0 & 19,0 & 12,7         & 1,72 & 3,39 & 6,75 & 3,95 \\
    \midrule
    \multirow{3}{*}{D-Fire}
      & E-CapsNet & 2,9  & 4,6  & 5,4  & 4,3          & 1,44 & 1,62 & 1,38 & \textbf{1,48} \\
      & ResNet-18 & 6,8  & 5,3  & 4,7  & 5,6          & 1,84 & 1,26 & 1,79 & 1,63 \\
      & DeiT-Tiny & 2,4  & 5,0  & 3,6  & \textbf{3,7} & 2,14 & 7,57 & 7,70 & 5,80 \\
    \bottomrule
  \end{tabular}
\end{table}

O Efficient-CapsNet tem $\Delta_f$ médio cerca de $1$~pp menor que o da
ResNet-18 nos dois conjuntos (Tabela~\ref{tab:res-dataeff}). A direção é a
prevista por H2, mas a diferença não é significativa ($p_{\text{Holm}} = 0{,}27$)
e não altera o ordenamento absoluto: no CIFAR-10 com augmentação forte, a
ResNet-18 com um terço dos dados já supera o modelo capsular com todos eles. Pelo
critério de dados para um alvo, a vantagem aparece no D-Fire sem augmentação: o
Efficient-CapsNet atinge $87{,}7\%$ com cerca de $40\%$ dos dados, e a ResNet-18
precisa de cerca de $59\%$.

O resultado mais nítido a favor do modelo capsular não é sobre robustez nem
acurácia: ele é a arquitetura menos sensível à semente nos dois conjuntos, e
essa estabilidade quase não depende do regime de augmentação. O DeiT-Tiny, ao
contrário, chega a $19{,}4$~pp de amplitude em uma única célula. Além disso, a
norma das cápsulas de saída, usada como confiança, separa acertos de erros em
dados limpos tão bem quanto o \emph{softmax} (AUROC de até $0{,}856$ no CIFAR-10
e $0{,}837$ no D-Fire), o que é uma propriedade de discriminação e não de
calibração~\cite{guo2017calibration}. Sob deslocamento essa vantagem desaparece.

\section{Veredito sobre as hipóteses}
\label{sec:res:verdict}

\begin{table}[htb]
  \centering
  \caption{Avaliação das hipóteses, com valores de $p$ corrigidos por Holm.}
  \label{tab:res-verdict}
  \small
  \begin{tabular}{l p{0.82\textwidth}}
    \toprule
    Hipótese & Veredito \\
    \midrule
    H1 & \textbf{Não sustentada.} A retenção sob rotação não supera a da
    ResNet-18 ($+1{,}1$~pp no CIFAR-10, $p_{\text{Holm}} = 0{,}22$; $-1{,}5$~pp
    no D-Fire, $p_{\text{Holm}} = 0{,}065$). A evidência contrária vem da
    acurácia absoluta, em que a ResNet-18 lidera todos os regimes com
    augmentação no CIFAR-10 por até $8{,}6$~pp. \\
    \addlinespace
    H2 & \textbf{Não sustentada.} $\Delta_f$ é cerca de $1$~pp menor para as
    cápsulas nos dois conjuntos, sem significância ($p_{\text{Holm}} = 0{,}27$).
    Com três sementes, não se pode afirmar que H2 seja falsa. \\
    \addlinespace
    H3 & \textbf{Sustentada.} Os quatro efeitos encolhem sob augmentação forte:
    a diferença de retenção contra a ResNet-18 vai de $+5{,}6$ a $-2{,}0$~pp no
    CIFAR-10 e de $-3{,}5$ a $+1{,}4$~pp no D-Fire; a de $\Delta_f$, de $-2{,}2$
    a $0{,}0$~pp e de $-3{,}9$ a $+0{,}7$~pp. \\
    \bottomrule
  \end{tabular}
\end{table}

O achado central é negativo e vale para este delineamento: substituir camadas
convolucionais por cápsulas com roteamento por autoatenção não produziu ganho
líquido de robustez, e o ganho de eficiência de dados é menor que a perda de
acurácia em distribuição que o acompanha (Tabela~\ref{tab:res-verdict}). Há,
porém, uma qualificação que os números impõem. Sem augmentação, o
Efficient-CapsNet retém $5{,}6$~pp a mais de acurácia sob rotação que a
ResNet-18 no CIFAR-10. A vantagem descrita pela literatura existe quando as
linhas de base não recebem augmentação, e desaparece quando recebem. O resultado
não é que as cápsulas nunca tiveram a propriedade, e sim que ela deixou de ser
distintiva.